\section{Closure matrix by load-bearing theorem node}\label{sec:p9-matrix}

\begin{description}[leftmargin=3.3cm,style=nextline]
\item[Core Theorem~I.1 --- passed] Exact endpoint wording, qualitative faithful-Wilson universality, and the explicit non-claims are frozen against overreading. \emph{Witness:} Core Sections~I and~VII; proof-kernel Sections~1 and~7.

\item[Core Section~II packet map --- passed] Packet~8 now visibly consumes the Lane~A seam nodes and the public spine no longer hides that dependence. \emph{Witness:} Core Section~II and Appendix~A; proof-kernel dependency-spine section.

\item[IV.2 / III.18 --- passed] Same-scale Wilson return is synchronized everywhere at finite logarithmic depth rather than vague fixed depth. \emph{Witness:} Core Section~IV; Companion~I theorem-local note at III.18.

\item[V.6 / 5.69 --- passed] Finite-truncation inverse control carries an explicit downstream-use ledger into the finite-cap bridge and closure packet. \emph{Witness:} Core Section~V; Companion~II theorem-local note at 5.69.

\item[V.9 / F.331 --- passed] The mixed-correlator closure node no longer floats in isolation; it now feeds a theorem-facing bridge into 5.74. \emph{Witness:} Core Section~V; Companion~II state-construction ledger and bridge notes.

\item[V.10 / 5.74 --- passed] Finite-cap state existence, positivity, unitality, and OS structure are theorem-faced locally instead of reconstructed from distributed hints. \emph{Witness:} Core Section~V; Companion~II theorem 5.74 together with the theorem-facing bridge notes.

\item[V.11 / 5.76 --- passed] The inductive-union passage remains sharply separated from 5.75 compatibility language across all proof homes and summaries. \emph{Witness:} Core Section~V; Companion~II theorem-local note at 5.76; freeze appendix crosswalk.

\item[V.13 / 5.77 --- passed] Bounded-base cyclicity remains a separate theorem-facing output that is consumed transparently at QE3 and reconstruction. \emph{Witness:} Core Section~V; Companion~II theorem-local note at 5.77.

\item[F.212 / VI.2 --- passed] The QE3 seam theorem now cites the correct nodes with the explicit local split 5.75 / 5.74(b) / 5.76 / F.213 / F.214. \emph{Witness:} Core Section~VI; Companion~I theorem F.212 and its local note.

\item[F.216 / VI.7 --- passed] Continuum vacuum-gap transport keeps the density / graph-core handoff explicit and does not hide a domain argument in prose. \emph{Witness:} Core Section~VI; Companion~I theorem-local note at F.216.

\item[M.1 / VII.1 --- passed] The Euclidean OS dossier remains the exact internal endpoint handoff, with the imported standard background isolated rather than smeared. \emph{Witness:} Core Section~VII; Companion~III theorem M.1; proof-kernel imported-background appendix.

\item[O.3 / VII endpoint --- passed] Qualitative faithful-Wilson universality remains visibly distinct from the $\rho$-indexed quantitative non-claim and from any nonlocal-sector overreading. \emph{Witness:} Core Section~VII; Companion~III endpoint section; proof-kernel Section~VII.
\end{description}

The purpose of this matrix is not to restate every proof. It is to record that each load-bearing node singled out in the closure plan now has a theorem-facing witness inside the packaged suite and therefore no longer needs reconstructive labor to see why it sits on the public spine.
